\chapter{Guix System: One File to Rule Them All}
\label{chap:guix-system}

\epigraph{``A configuration management system like Ansible or Puppet is like a robot that walks around your kitchen moving spices. Guix System is a 3D printer that manufactures a pristine kitchen from scratch on every change.''}{--- Overheard at FOSDEM}

\noindent
Up to this point, we have used Guix on top of a foreign distribution (like Debian, Fedora, or Arch). But Guix is not just a package manager; it is also a complete, standalone operating system: \textbf{Guix System}.

In Guix System, your entire machine---kernel, init system, system daemons, user accounts, filesystems, network configuration, and desktop environment---is declared inside a single Guile Scheme file, typically called \texttt{/etc/config.scm}.

\section{The Imperative OS Trap vs The Declarative OS}

On an ordinary Linux system:
\begin{itemize}
  \item You edit \texttt{/etc/network/interfaces} or fiddle with NetworkManager.
  \item You manually edit \texttt{/etc/fstab} and pray you didn't typo the UUID.
  \item You run \texttt{useradd}, \texttt{usermod}, and \texttt{visudo}.
  \item You install PostgreSQL, tweak \texttt{postgresql.conf}, and run \texttt{systemctl enable postgresql}.
\end{itemize}

Six months later, your server is a delicate, undocumented snowflake. Nobody knows which packages were installed manually, which configuration files were edited, or how to recreate the machine if the hardware dies.

In Guix System, \textbf{the file is the machine}. If a service or package is not declared in your Scheme configuration, it does not exist on your system.

\subsection{The Imperative Temptation vs. Declarative Reproducibility}

Having read about package management in Chapter~\ref{chap:cli-daily-life}, your very first instinct after booting into Guix System might be to open a shell and type:
\begin{center}
\guixcmd{guix install git emacs firefox htop}
\end{center}

\textbf{Do not do this!} This is the single most common pitfall for newcomers arriving from traditional distributions like Debian, Arch, or Fedora.

In Chapter~\ref{chap:cli-daily-life}, we emphasized that the non-declarative, imperative workflow (\guixcmd{guix install}, \guixcmd{guix remove}) is primarily intended for users running Guix as a supplementary package manager on top of a \textbf{foreign distribution}. On Guix System, relying on \guixcmd{guix install} actively destroys one of the operating system's greatest superpowers: \textbf{complete, mathematical reproducibility}.

Consider what happens if you install software imperatively on Guix System:
\begin{itemize}
  \item Packages are placed in an isolated, untracked profile under \texttt{\textasciitilde/.guix-profile}.
  \item Those packages remain completely invisible to your declarative operating system definition.
  \item If your SSD fails tomorrow, or you wish to provision an identical backup laptop or cloud server from your Git repository, those imperatively installed packages will be completely absent because they were never recorded in your configuration code!
\end{itemize}

On Guix System, \textbf{the official, recommended, and reproducible way to install new packages is declaratively through \texttt{/etc/config.scm}}. Instead of mutating the machine state with ad-hoc terminal commands, you declare your desired software in Scheme code, commit it to version control, and instantiate it with \guixcmd{guix system reconfigure}.

\section{The Anatomy of \texttt{config.scm}}

Let us examine a complete, fully functional \texttt{config.scm} for a desktop workstation:

\begin{lstlisting}[style=guixscheme]
;; /etc/config.scm - Declarative Guix System configuration
(use-modules (gnu)
             (gnu system nss))
(use-service-modules desktop networking ssh sddm)
(use-package-modules certs fonts gnome linux package-management)

(operating-system
  (host-name "guix-battlestation")
  (timezone "Europe/Paris")
  (locale "en_US.utf8")

  ;; Bootloader Configuration (UEFI GRUB)
  (bootloader (bootloader-configuration
                (bootloader grub-efi-bootloader)
                (targets '("/boot/efi"))
                (keyboard-layout keyboard-layout)))

  ;; Storage & Filesystems
  (file-systems (cons* (file-system
                         (mount-point "/")
                         (device (file-system-label "my-root"))
                         (type "ext4"))
                       (file-system
                         (mount-point "/boot/efi")
                         (device (uuid "1234-ABCD" 'fat))
                         (type "vfat"))
                       %base-file-systems))

  ;; User Accounts
  (users (cons (user-account
                 (name "alice")
                 (group "users")
                 (supplementary-groups '("wheel" "netdev" "audio" "video" "kvm"))
                 (home-directory "/home/alice"))
               %base-user-accounts))

  ;; System-wide packages
  (packages (append (list nss-certs     ; HTTPS Certificates
                          font-dejavu   ; Core fonts
                          git
                          emacs)
                    %base-packages))

  ;; System Daemons and Desktop Services
  (services (append (list (service openssh-service-type
                            (openssh-configuration
                              (port-number 2222)
                              (permit-root-login #f)))
                          (service gnome-desktop-service-type))
                    %desktop-services)))
\end{lstlisting}

\subsection{Installing New Packages: The Declarative Workflow}

Examine lines 67--71 in the configuration above. This \texttt{packages} field is where system-wide software installation actually happens on Guix System:

\begin{lstlisting}[style=guixscheme]
  ;; System-wide packages
  (packages (append (list nss-certs     ; HTTPS Certificates
                          font-dejavu   ; Core fonts
                          git
                          emacs)
                    %base-packages))
\end{lstlisting}

When you want to install a new package on your machine, you do not execute an imperative terminal command. Instead, you follow a clean four-step declarative workflow:

\begin{enumerate}
  \item \textbf{Search}: Find the package name using \guixcmd{guix search <query>}.
  \item \textbf{Import}: Ensure the module containing the package definition is imported at the top of your \texttt{config.scm} via \texttt{use-package-modules} (for example, \texttt{(use-package-modules admin)} for \texttt{htop}, or \texttt{(use-package-modules version-control)} for \texttt{git}).
  \item \textbf{Declare}: Add the package variable name to the \texttt{(list ...)} inside the \texttt{packages} field.
  \item \textbf{Reconfigure}: Apply your changes across the operating system:
\begin{lstlisting}[style=guixcli]
sudo guix system reconfigure /etc/config.scm
\end{lstlisting}
\end{enumerate}

Guix immediately downloads pre-built substitutes or builds the package, creates a brand-new operating system generation, and links the binaries into the system profile. If you later decide to uninstall the software, you simply delete its line from \texttt{config.scm} and reconfigure again.

\begin{alchemybox}[Why Declarative Code Trumps Imperative Commands]
By declaring packages directly inside \texttt{config.scm}, you guarantee that:
\begin{itemize}
  \item \textbf{The Machine Is Source Code}: Your entire software environment is captured in plain text. You can commit your configuration to Git, review software additions in pull requests, and audit changes over years.
  \item \textbf{Effortless Replication}: Setting up a second workstation or replacement laptop takes zero manual guesswork. Clone your configuration, run \guixcmd{guix system reconfigure}, and you get an identical system down to the last library.
  \item \textbf{Synchronized Rollbacks}: If an added package introduces a bug or breaks your workflow, rolling back to the previous system generation in GRUB reverts the packages, kernel, and system daemons simultaneously as one cohesive unit.
\end{itemize}
\end{alchemybox}

\begin{alchemybox}[Service Composition with \texttt{append} and \texttt{modify-services}]
Notice how services are combined using standard Scheme list operations!
\texttt{\%desktop-services} is simply a Scheme list containing dozens of core services (udev, elogind, dbus, network-manager, polkit, etc.).
You can add services with \texttt{append} or \texttt{cons}, or modify existing service configurations using the \texttt{modify-services} macro:
\begin{lstlisting}[style=guixscheme]
(modify-services %desktop-services
  (guix-service-type config =>
    (guix-configuration
      (inherit config)
      (substitute-urls
        '("https://ci.guix.gnu.org"
          "https://bordeaux.guix.gnu.org")))))
\end{lstlisting}
\end{alchemybox}

\section{Real-World Hardware: Using Nonguix for Laptops \& Wi-Fi}

By default, Guix System uses the 100\% free \textbf{Linux-Libre} kernel. If you are running on modern ThinkPads, MacBooks, or gaming laptops with proprietary Intel/Realtek Wi-Fi and Bluetooth chips, you can use the \textbf{Nonguix} channel to declare the standard non-libre Linux kernel and firmware:

\begin{lstlisting}[style=guixscheme]
;; Laptop configuration with Nonguix kernel and firmware
(use-modules (gnu)
             (nongnu packages linux)     ; From nonguix channel
             (nongnu system linux-initrd))

(operating-system
  (inherit %base-operating-system)
  (host-name "nomad-laptop")

  ;; Use vanilla Linux kernel with proprietary binary blob drivers
  (kernel linux)
  (initrd microcode-initrd)
  (firmware (list linux-firmware))

  ;; Automatically authorize the Nonguix pre-built substitute server!
  (services
    (modify-services %desktop-services
      (guix-service-type config =>
        (guix-configuration
          (inherit config)
          (substitute-urls
            (append (list "https://substitutes.nonguix.org")
                    %default-substitute-urls))
          (authorized-keys
            (append (list (local-file "/etc/nonguix-key.pub"))
                    %default-authorized-guix-keys)))))))
\end{lstlisting}

With this declaration, your laptop gets full hardware acceleration, functional Wi-Fi, and CPU microcode security patches without losing any of Guix System's declarative purity or atomic rollback powers!

\section{Applying the System: \texttt{guix system reconfigure}}

To apply your changes, you execute:

\begin{lstlisting}[style=guixcli]
sudo guix system reconfigure /etc/config.scm
\end{lstlisting}

When you run this command:
\begin{enumerate}
  \item Guix instantiates and validates all Scheme definitions.
  \item It builds all required store paths in \texttt{/gnu/store}.
  \item It builds an entirely new operating system generation.
  \item It signals the Shepherd init system to reload modified daemon configurations, spawn new services, and stop retired services.
  \item It updates the GRUB bootloader menu with a new entry: \textbf{Generation $N+1$}.
\end{enumerate}

\section{The Unbreakable Bootloader: GRUB Generations}

What happens if you accidentally misconfigure your X server or build a kernel that kernel-panics on boot?

On other distributions, you are frantically reaching for a live USB stick and mounting partitions via \texttt{chroot}.

On Guix System, you simply reboot your machine. When the GRUB boot menu appears, you select:

\begin{center}
\textbf{GNU Guix System --- Generation 42 (Yesterday at 18:30)}
\end{center}

You hit \texttt{Enter}, and the system boots into the exact, known-working state of yesterday's operating system. 

\begin{tearsbox}[The Broken Kernel Nightmare]
Every veteran Linux user has an emotional scar from a bad kernel upgrade breaking their graphics drivers or Wi-Fi. In Guix System, an operating system update is never destructive. Every past system generation remains independently bootable in GRUB until you explicitly choose to garbage-collect it.
\end{tearsbox}

\begin{figure}[htbp]
\centering
\includegraphics[width=0.92\textwidth]{images/rollback_button.jpg}
\caption*{\textit{Figure 6.1: The ultimate sysadmin superpower: When an experimental update catches fire, casually tap the GRUB Rollback button to rewind reality.}}
\end{figure}

\section{Building Images: ISOs, VMs, and Cloud AMIs}

Because your system is a pure Scheme function, Guix can transform your \texttt{config.scm} into any output format with a single command:

\begin{lstlisting}[style=guixcli]
# 1. Build a bootable installation ISO image
guix system image -t iso9660 /etc/config.scm

# 2. Build a QEMU virtual machine image (QCOW2)
guix system image -t qcow2 /etc/config.scm

# 3. Build a Docker container image representing the complete OS
guix system image -t docker /etc/config.scm
\end{lstlisting}
